\chapter{Introduction}
\label{ch:introduction}

The standard model of cosmology describes a universe that emerged from
a hot, dense initial state and has been expanding and cooling ever
since. As it cooled, the fundamental forces of nature underwent a
series of symmetry breaking transitions, each leaving behind
observable imprints on the structure of the universe we see today.
Understanding these transitions from first principles requires a
synthesis of quantum field theory, statistical mechanics, and
cosmology, and the central tool that makes this synthesis possible is
the effective action.

This thesis develops that tool systematically, beginning with its
abstract foundations and ending with concrete physical predictions for
gravitational wave signals detectable by next-generation experiments.
The thread connecting all five parts is a single question: how does
the vacuum structure of a quantum field theory change as a function of
temperature, and what are the cosmological consequences of those
changes?

\section*{The Effective Action and the Quantum Vacuum}

In classical field theory the vacuum is simply the minimum of the
classical potential. In quantum field theory this picture is
corrected by loop diagrams, which shift the location of the minimum,
change the shape of the potential, and can even generate new minima
that have no classical counterpart. The object that encodes all of
these quantum corrections is the effective action $\Gamma[\varphi_c]$,
the generating functional for one-particle-irreducible Green's
functions. Its equation of motion,
$\delta\Gamma/\delta\varphi_c = 0$, governs the full quantum dynamics
of the field with all loop corrections included, and its restriction
to constant field configurations defines the effective potential
$V_{\text{eff}}(\varphi_c)$, whose minima are the true vacua of the
quantum theory.

The construction of the effective action from the path integral,
through the generating functionals $Z[J]$ and $W[J]$ and the Legendre
transform, is developed in Part~I. The physical content of this object
is richer than its definition suggests. The effective potential is
simultaneously the generating functional for zero-momentum 1PI
diagrams, the minimum energy density over all states with a given
field expectation value, and the quantum-corrected potential that
appears in the equation of motion for the classical field. All three
interpretations are developed and used.

The explicit computation of the effective potential at one loop, first
in scalar field theory and then in gauge theories, reveals one of the
most striking consequences of quantum corrections: dimensional
transmutation. In massless scalar electrodynamics, a theory with no
classical mass scale and a classically scale-invariant action, quantum
corrections generate a nontrivial minimum of the effective potential
at a nonzero field value, dynamically producing a mass scale from
a dimensionless coupling constant. This mechanism, due to Coleman and
Weinberg, is the prototype for a broad class of phenomena in which
the vacuum structure of a theory is qualitatively different from what
classical analysis would suggest. It also provides the first example
of a theme that recurs throughout the thesis: quantum corrections do
not merely shift numbers, they can change the qualitative character
of the vacuum.

\section*{Finite Temperature and Phase Transitions}

The effective potential computed in Part~I describes the vacuum of the
quantum field theory at zero temperature. In the early universe, matter
was in thermal equilibrium at high temperatures, and the relevant
object is not the ground state energy but the free energy. Part~II
extends the effective action formalism to finite temperature using the
imaginary-time formalism of thermal field theory, in which the inverse
temperature $\beta = 1/T$ plays the role of a compactified Euclidean
time direction and quantum fields satisfy periodic or antiperiodic
boundary conditions.

The key result of this extension is that finite-temperature
corrections generate a positive mass term proportional to $T^2\varphi_c^2$,
which competes with any negative mass term in the zero-temperature
potential. This competition drives symmetry restoration: even if the
zero-temperature theory has a spontaneously broken ground state, above
a critical temperature $T_c$ the symmetric phase becomes the global
minimum of the finite-temperature effective potential. As the universe
cooled through $T_c$, the symmetry was broken and a phase transition
occurred. The order and character of that transition, whether it
proceeds smoothly or via bubble nucleation, depends on the detailed
structure of the effective potential and has profound observational
consequences.

Part~III applies this framework to the cosmological setting. The
thermal history of the early universe is punctuated by a sequence of
such symmetry breaking transitions: a possible GUT transition at
$T\sim 10^{16}$ GeV, the electroweak transition at $T\sim 100$ GeV,
and the QCD confinement transition at $T\sim 200$ MeV. For each of
these, the effective potential, computed using the tools developed in
Parts~I and~II, determines whether the transition is first or second
order and what the thermodynamic parameters of the transition are.

\section*{False Vacuum Decay and Bubble Nucleation}

When the effective potential has two minima separated by a barrier,
the universe can find itself temporarily trapped in the metastable
false vacuum. The transition to the true vacuum does not proceed by
the field rolling classically over the barrier but by quantum and
thermal fluctuations nucleating bubbles of the true vacuum within
the surrounding false vacuum. Each bubble expands, driven by the
pressure difference between the two phases, and eventually the bubbles
percolate, completing the transition.

The quantitative description of this process, developed in Part~IV,
is one of the most elegant applications of the path integral formalism
in quantum field theory. The nucleation rate is controlled by the
bounce solution, a saddle-point of the Euclidean action that
interpolates between the false and true vacua. At zero temperature
the relevant bounce has $O(4)$ symmetry, while at finite temperature
it has $O(3)$ symmetry, reflecting the compactification of Euclidean
time. In both cases the nucleation rate takes the form
$\Gamma \sim e^{-B}$, where $B$ is the Euclidean action of the bounce,
and the dynamics of the resulting bubbles are determined by analytic
continuation of the bounce solution from Euclidean to Minkowski space.

The phase transition is characterized by three dimensionless
parameters: $\alpha$, the ratio of the latent heat to the radiation
energy density; $\tilde\beta/H$, the ratio of the transition rate to
the Hubble rate; and $v_w$, the bubble wall velocity. These parameters
determine the gravitational wave spectrum produced during the
transition through the three mechanisms of bubble wall collisions,
sound waves in the primordial plasma, and magnetohydrodynamic
turbulence. For an electroweak-scale transition, the predicted
spectrum peaks at millihertz frequencies, precisely the sensitivity
band of the LISA space interferometer, providing a direct
observational window into the thermal history of the early universe
at energies far beyond the reach of terrestrial accelerators. The key
references for this gravitational wave signal are Caprini et al.
(JCAP 2016, arXiv:1512.06239) and Caprini et al. (JCAP 2020,
arXiv:1910.13125).

\section*{Quantum Fields in Curved Spacetime and Semiclassical Gravity}

All of the above is developed in flat spacetime, or at most in a
classical FRW background treated as fixed. In the early universe,
however, the curvature of spacetime is not negligible, and a more
complete treatment is required. The appropriate framework is
semiclassical gravity, in which the spacetime geometry is treated
classically via the Einstein equations but the matter fields are
promoted to quantum operators. The backreaction of quantum matter on
the geometry is captured by the semiclassical Einstein equation,
\begin{equation*}
    G_{\mu\nu} = 8\pi G\langle\hat T_{\mu\nu}\rangle,
\end{equation*}
where the right-hand side is the expectation value of the quantum
stress-energy tensor in some state. This is the natural next step
beyond the flat-spacetime effective action formalism developed in the
earlier parts of the thesis, and it is valid as long as the curvature
remains well below the Planck scale so that quantum gravitational
effects themselves can be ignored.

The necessity of this framework is already visible in the effective
action. In a curved background the field equation acquires a
non-minimal coupling to the curvature,
$(\square - m^2 + \xi R)\phi = 0$, which shifts the effective mass
and can qualitatively change the phase structure of the theory.
More fundamentally, semiclassical gravity reveals that the vacuum
itself is observer-dependent. In flat spacetime, the Poincar\'e
invariance of the Minkowski metric singles out a unique vacuum state
agreed upon by all inertial observers, and the effective potential
is a well-defined object. In a general curved spacetime there is no
such preferred frame. Different observers define inequivalent mode
decompositions of the quantum field, related by Bogoliubov
transformations, and the vacuum of one observer contains particles
as seen by another.

Part~V develops this formalism systematically. We begin with the
causal structure of the Schwarzschild geometry, introducing the
Kruskal and Penrose diagrams that make the global structure of
spacetime transparent. Black hole thermodynamics follows naturally:
the area of the horizon never decreases, in direct analogy with
entropy, and a gedanken experiment involving a qubit lowered into a
black hole fixes the Bekenstein-Hawking entropy $S_{BH} = A/4$ and
the generalized second law $\delta S_{\text{tot}} \geq 0$. The
general Bogoliubov transformation formalism is then developed,
leading to the key result that the vacuum of one observer can
contain a thermal distribution of particles as seen by another
whenever the $\beta$ coefficients are nonzero.

The physical consequences are explored through two concrete examples
that build toward Hawking radiation. The moving mirror shows that
even in flat spacetime a time-dependent boundary condition produces
a thermal flux of particles when the mirror asymptotically approaches
a null ray, and the logarithmic relation between the in and out
coordinates near that asymptote is the essential mathematical
structure. The Unruh effect shows that a uniformly accelerating
observer in the Minkowski vacuum measures a thermal bath at
temperature $T = a/2\pi$, where $a$ is the proper acceleration,
because the causal horizon of the accelerating observer plays
precisely the same role as the horizon of a black hole.

These two results are not analogies. They are the same physics as
Hawking radiation, seen in simpler settings. When a black hole forms
by gravitational collapse, the logarithmic relation between the
coordinates at past and future null infinity produces, in the
semiclassical gravity framework, a steady thermal flux of radiation
at the Hawking temperature $T_H = \kappa/2\pi$, where $\kappa$ is
the surface gravity of the horizon. The black hole emits as if it
were a black body at this temperature, and the resulting evaporation
has deep consequences for the consistency of quantum mechanics and
general relativity that remain incompletely understood.

The connection back to the rest of the thesis is direct. The
de~Sitter spacetime produced during inflation also possesses a
cosmological horizon, and quantum fields in that background
experience a Gibbons-Hawking temperature $T = H/2\pi$ by the same
mechanism. Thermal fluctuations in this background seed the
primordial density perturbations that grow into the large-scale
structure of the universe, completing the circle from the microscopic
effective action of Part~I to the observable universe of today.\section*{Outline}

The thesis is organized as follows. Part~I develops the effective
action formalism at zero temperature, computing the one-loop effective
potential in scalar theory, massless scalar electrodynamics, the
Glashow-Salam-Weinberg electroweak model, and non-Abelian gauge
theories, with renormalization carried out both by hard cutoff and by
dimensional regularization in the $\overline{\text{MS}}$ scheme.
Part~II develops thermal field theory using the imaginary-time
formalism and computes finite-temperature corrections to the effective
potential for scalar fields, the GSW model, and minimal SU(5).
Part~III places these results in their cosmological context, covering
the thermal history of the early universe, the characterization of
first-order phase transitions, and the gravitational wave signatures
they produce. Part~IV analyzes false vacuum decay and bubble
nucleation via the bounce solution, at both zero and finite
temperature. Part~V treats quantum field theory in curved spacetime,
from the causal structure of the Schwarzschild geometry through black
hole thermodynamics, the Bogoliubov transformation formalism, and the
Unruh effect, culminating in the derivation of Hawking radiation and
its physical interpretation. Technical derivations are collected in
the appendices.
